\chapter{Information, Causation, and the Blind Spot}
\label{chap:pledge-ch2}

\section{Information and Causation}
\label{sec:pledge-iv}

In the early 2000's i attended a workshop at the W3C ostensibly about developing an ontology for supporting online communication about biology. At the time i was already aware of David Harel's work suggesting that genes were not at all the right way to organize the picture of how heritable protein production actually happens. As i sat in the workshop observing the gap in consensus and the rate at which consensus was emerging i felt a growing sense that this effort was more doomed than the Sisyphean task of painting the Golden Gate Bridge. You see by the time they complete painting the bridge it is time to begin painting it again.

i was keenly aware of a sense that by the time any sort of ontology had achieved acceptance within the body of organizations assembled at the workshop the proposal would not only be out of date, but that cracks in its very foundations would be showing. In fact, cracks were already showing and they had barely begun . It left a profound impression on me. i began to wonder about the very enterprise of classification. Are all classification schemes grass: doomed to spring up in multitudes and be mowed down by subsequent experiments?

To be fair some classification schemes seem more robust. For example, during the previous decade i witnessed the extraordinary work on the classification of all finite simple groups. The very building blocks of symmetry had been carefully teased apart and classified with compelling proofs buttressing the construction. i wondered whether there was always going to be a gap between the clean lines of mathematics and the significantly blurrier lines of the sciences. Could there be a sense in which these two disciplines informed each other about the very nature of classification?

While contemplating these questions i was also engaged in a small side project about employing computation as a source of invariants. In classical mathematics algebraic entities, like groups, are used to record information about other entities, such as topological spaces. For example, mathematicians had learned that you could glean essential information about the nature of a space by studying the collections of loops you could trace in the space. If the space has a hole in it, as a donut or a coffee cup does, then you can't contract the loops that circumnavigate the hole down to a point. The different kinds of loops form the building blocks of a group that can be used to classify the space. This basic idea --- using algebraic information to classify other kinds of entities --- has many offshoots and has had profound implications even for the foundations of mathematics.

i observed that algebraic entities don't enjoy the same sorts of dynamics that computations do. While a group might be a representation of dynamics, the abstraction itself just sits there. It doesn't wiggle and squirm the way a computation does. Dynamics lives in the maps between the algebraic entities. This way of looking at dynamics is in fact codified in the category theoretic account of classical mathematics. Mathematicians study the category of groups and group homomorphisms, or the category of topological spaces and homeomorphisms. But computation carves up the world differently. 

This is much easier to see in the computational calculi like the $\lambda$-calculus or the $\pi$-calculus or the rho calculus, but i want to defer a more formal investigation a bit longer. Fortunately, it's still fairly obvious in more ad hoc constructions like the Turing machine. Although we have to think about things like the configuration of the machine, such as what is on the tape at the beginning, given such a configuration we intuitively understand that the computation proceeds autonomously. We don't necessarily have to engage it. This way of thinking about things, this halfway house to agency, puts the dynamics right on the structure that represents the computation. \emph{It doesn't put the dynamics on maps between Turing machines}. Or even on maps between collections of Turing machines. 

This is a big shift. Indeed, at the time there was no widely accepted category theoretic account of computation that illuminated this fundamental difference between how classical mathematics carves up the world and how computation carves up the world. When we get to the formal presentation we will present a category in which the objects are models of computation and the maps between them preserve something called bisimulation. The objects represent the dynamics and the morphisms between them preserve these dynamics. But again i'm getting ahead of myself.

The key idea i was exploring was to associate computational entities --- something like Turing machines --- to topological entities, effectively retracing the early ideas in homotopy theory, but substituting computations for algebraic structures. My first foray into this was associating to each knot a term in the $\pi$-calculus. This has some interesting consequences and deserves a much deeper dive than we can go into here. The relevant point here is that it set me down the path of thinking about classification in terms of bisimulation. There's that word again. What is bisimulation?

We are going to have a lot to say about bisimulation in the Turn. But there's a very intuitive way to understand the concept that is encoded in memetic culture. The phrase ``anything Bob can do Alice can do'' is half the notion. Paired with ``and anything Alice can do Bob can do'' we complete the picture. One of the many interesting aspects of this formulation is its can-do attitude. What does it mean to say Alice or Bob can do something? In subsequent discussion it will become clear that ``can do'' is dual to experiment. Distinguishing between two agents on the basis of what they can do is equivalent to distinguishing them in terms of experiments. 

This approach has profound implications for classification. Everything is sorted into what can be distinguished by experiment. If there is no experiment that distinguishes between two things, they are the same. Classification ceases to be a Sisyphean task. We can move from asking grok to give us a break down of the news to actually grokking a situation, circumstance, idea, or structure. At least when we are talking about computation. If the world is something other than a computation, then all bets are off!

\section{Measurement}
\label{sec:pledge-v}

We interrupt your regularly scheduled reading with this brief detour. In 2021 i proposed a collection of games that humans are guaranteed to win against machines, at least for now. One of them i called speed breath chess. The main difference between this and speed chess is that the player has to hold their breath when it is their turn to move. When it is the machine's turn they lose, because they don't breathe. Further, if we use electrical power as a proxy for breath, then when it is the machine's turn we turn off their power. Again they lose. The point is that access to breath is part of what makes us human and connects us to all living things.

i have been contemplating access to breath quite a lot. My partner, the one i mentioned with whom i have found a qualitatively different relationship, is now in late stage ALS. In the last 6 weeks her access to breath has been dramatically compromised as her lungs have weakened precipitously. i had an episode in my life that i thought gave me a sense of the challenges she is facing. i thought i understood the kind of fear that arises in the body, quite apart from any narrative the mind might spin, or the heart might score the soundtrack to.

In the early 90's i moved to London with my family to pursue a PhD at Imperial College under Samson Abramsky. The summer of our first year in London there was a rare confluence of events in one week that left me literally breathless. First, there was a historic surge in ragweed pollen --- which i didn't know i was allergic to at the time. Second, there was a thermal inversion over the city, trapping the pollution. Third, there was a train strike, causing record numbers of Londoners to use their automotives. i remember leaving the house quite early to get a swim in before assuming my academic duties. i was in excellent shape and regularly swam a mile in the morning, and later biked, and ran. By the time i left the pool that day i was very short of breath. On my way to class i knew something was badly wrong and headed home. Within an hour i was flat on my back and felt like i was running a 4 min mile. When we called the doctors they rushed out to administer breathing aid. Apparently many people died that day due to the unfortunate convergence of conditions.  

In the last few weeks i have called on this memory often as an aid in understanding how my partner must be feeling. But, last night we had a conversation that completely changed my sense of things. Because her speech is now dramatically compromised, she held up her end of the conversation by spelling her thoughts out on a tablet with her eyes. During the course of our conversation i decided to probe my understanding of her experience of breath and asked her about it. As she described her initial terror in the situation, followed by an interim stabilization, something dawned on me. 

In the previous months she had been using a device called a bipap (relative of a cpap) to support her breathing, but only intermittently. About an hour or two during the day and then throughout the night. But then one night in the transition from her recliner, where she spends most of her day, to her bed her breathing got so shallow she blacked out. Her caregivers called emergency medical services, and the paramedics came, but it was actually her caregivers who revived her with judicious use of the bipap. If i understand her, at that point a really deep anxiety set in --- one i thought i could understand through my experience in London. 

To stabilize her condition she shifted to continuous use of the bipap. She wears a mask all day and all night connected to a machine that helps her breathe. To help make the anxiety manageable her healthcare provider recommended a regimen of morphine. Dear Reader, if you are not aware, this marks a major shift in care, from maintenance and management to palliative care. Morphine relieves the anxiety at the cost of depressing the respiratory system. The regime becomes one of narrowing gyres towards the center of a mystery we each face. 

As she described her experience it landed in me quite differently. My partner is a profoundly spiritual person, with a personal practice developed through many years of practical application as a Buddhist, a school counselor, and a good friend to many. One of the points of connection between us is that both of us have come to a relationship with breath that is simultaneously practical and a doorway. It's more than a mere narrative about breath as the root of aspiration and inspiration, or about the etymological connection between breath and spirit. For myself, when people ask me about it, i say that this doorway is the one through which whatever i'm seeking can come find me. It took me decades to find this relationship to breath and i rely on it. It helps me self-regulate, to keep skyrocketing levels of cortisol balanced out, when the stresses of life threaten to overwhelm. And it is one of my truest delights is when my partner and i can just breathe together.

As we spoke something i understood intellectually became a felt experience. That relationship with breath, with spirit, is a dynamic one. Like my partner, one day i will feel the breath begin its departure from the body. That rhythm, that cadence, that constant companion whose undulations trace a silver thread between the mystery and embodiment will one day loosen its hold of me. This dynamism and impermanence --- this hazard --- is what makes the relationship to breath, to spirit, really alive, really worth deepening.

So i ask you, Dear Reader, if the mind we find is devoid of a connection to breath, to spirit, to this hazardous cadence, is it the kind of mind we really want to find? As we spend billions of dollars and kilowatt hours and mega liters of water on this doorway we call AI, is this the mind we want to come find us?

\section{The Blind Spot}
\label{sec:pledge-vi}

Part of the reason for the apparent non sequitur of the last section was to begin to grapple with the stakes of this question. There are consequences both to our outer and inner experience. As i mentioned in a previous article, the last time Life introduced an adaptability protocol like Mind the consequences were dire for all the other species on planet Earth, and the species hosting Mind is currently destroying their own ecological niche. If that species begets another kind of Mind, a speciation of Mind, if you will, but this time not directly connected to the substrate of Life, not directly connected to breathing or breeding, it's quite likely to recapitulate the previous emergence of Mind, but on an even grander scale. That is, to wipe out that which is not its own kind, and then wipe out its own means of sustaining itself. These are possible consequences for our outer experience.

This is not idle speculation. Look at how data centers are gobbling up land, electricity, and water. Seriously, look at your own electricity bill. There are real, felt consequences to humanity's outer experience that come from how it frames and answers this question. It's much like there were real, felt consequences to humanity's exploration of the questions of the nature of physical reality. Tiny little equations like E = Mc\textasciicircum{}2 gave rise to the Manhattan Project, the bombing of Hiroshima and Nagasaki, mutual assured destruction, and much more. The sort of Mind we look for will have consequences on this scale or greater, since it was Mind that came up with standard model physics.

And \ldots{} the point of the last section was to connect our exploration to the consequences to our inner experience, as well. Do we want to want to find a Mind that is disconnected to breath? It's not a rhetorical question. In his famous scifi trilogy (Out of the Silent Planet, Perelandra, and That Hideous Strength) C.S. Lewis describes the mode of angelic being as that which neither breathes nor breeds. Does the Mind hosted in humanity long to become angelic, in this sense? Or to beget an angelic form of Mind that is not subjected to the indignities of breathing and breeding? 

On one hand this is not possible in a physics that doesn't admit perpetual motion machines. There will be a replication protocol (breading), and there will be an energy exchange protocol (breathing). More on this later. On the other, it is worth noting that in the Abrahamic traditions, and others, humanity's position is exhalted because of its relationship to mortality. It would appear that in a number of the world's religious traditions humanity's super power is transformation, as opposed to coordination. Meanwhile, Lucifer's particular narrative notwithstanding, the angelic mode of being doesn't admit much in the way of transformation. 

To put a mathematician's spin on this, the ineffable infinite is by definition ubiquitous. Throw a stone in any direction and you hit the ineffable infinite. What's rare is the finite, or the compact. Put another way, imagine winding the real number line around like the grooves on a vinyl record and then hang that record on the wall, like a dart board. Now stand a few paces away and take a dart, the tip of which is so fine as to be a point (in the sense of the real line), and throw it at the dart board. Let's assume that your aim is good enough that you always hit the dart board. (This is not an onerous assumption because we can freely scale the dart board to whatever size we want, or alternatively you can stand as close to the board as you like.) What is the likelihood that you hit a counting number? It's 0. Even if you throw forever, you never hit a counting number. That's how rare the finite is.

From this perspective, the finite beings; the ones who have a beginning, a middle, and and end; the ones whe are privileged to bear mortality are the rare gems that the ineffable infinite might cherish, precisely because we are so damn hard to find. Ironically humanity, both at the phylogenic and the ontogenic level, finds the counting numbers first. Every human child learns to count (except possibly the Piraha?). And if you look at the history of humanity's understanding of quantity, we find the counting numbers before we find the uncountable ones.

But i digress. The point is to develop at least some appreciation of just how high the stakes of how we frame the question of the nature of Mind are, both in terms of our inner and outer experience. How we ask and go about answering this question will shape what it feels like on the inside and what it feels like on the outside. Both individually, and collectively. 

Note how natural it was to sort our experience into two categories, inner and outer. Likewise, how natural it is to talk about that in terms of two categories, individual and collective. This kind of sorting, or classification of phenomena is fundamental to agentic beings, such as ourselves. In the next section we will return to the nature and role of classification in agentic beings, i.e. beings equipped with the possibility of transformation.
